% =====================================================================
% main.tex  --  Selective SARP-TW paper skeleton
% Taylor & Francis Interact layout + NLM (numbered) reference style
% Target journal (primary): Transportmetrica B: Transport Dynamics
% Fallback journal: Journal of Intelligent Transportation Systems (JITS)
%   -> Both are T&F / interact.cls. To switch reference style to JITS
%      (author-year / APA), swap \bibliographystyle{tfnlm} for the APA
%      variant and adjust natbib options; body stays identical.
% =====================================================================
\documentclass[]{interact}

\usepackage{epstopdf}
\usepackage[caption=false]{subfig}
\usepackage{amsmath,amssymb,bm}
\usepackage{multirow}
\usepackage{algorithm}
\usepackage{algpseudocode}   % if unavailable, use algorithmic
\usepackage{url}

% ---- NLM numbered citation setup (do NOT edit; from T&F template) ----
\usepackage[numbers,sort&compress]{natbib}
\bibpunct[, ]{[}{]}{,}{n}{,}{,}
\renewcommand\bibfont{\fontsize{10}{12}\selectfont}
\makeatletter
\def\NAT@def@citea{\def\@citea{\NAT@separator}}
\makeatother

\theoremstyle{plain}
\newtheorem{theorem}{Theorem}[section]
\newtheorem{proposition}[theorem]{Proposition}
\theoremstyle{definition}
\newtheorem{definition}[theorem]{Definition}

% ---- author-defined macros (gather all here) ----
\newcommand{\Kmax}{K_{\max}}
\newcommand{\sarp}{\textsc{Selective SARP-TW}}

\begin{document}

\articletype{RESEARCH ARTICLE}

\title{A Reject-Aware Deep Reinforcement Learning Framework for
Selective Passenger--Freight Sharing in Demand-Responsive Off-Peak
Bus Scheduling with Hard Time Windows}

% TODO: before de-anonymized submission, replace with English author block:
%   \name{Pinyin Name\textsuperscript{a}\thanks{CONTACT ... Email: ...}}
%   \affil{\textsuperscript{a}English Department, English University, City, Country}
% NOTE: interact.cls compiles with pdfLaTeX and does NOT accept CJK characters.
% Keep the author block in English/pinyin only.
\author{
\name{Anonymous Author(s)}
\affil{Affiliation withheld for double-blind review}
}

\maketitle

% =====================================================================
\begin{abstract}
Off-peak demand-responsive transit (DRT) chronically under-utilizes vehicle
capacity, while urban logistics faces mounting last-mile pressure; co-loading
parcels alongside passengers is an attractive way to raise fleet efficiency,
provided that hard passenger service guarantees are preserved. We formalize the
off-peak ($10\!:\!00$--$16\!:\!00$) shared passenger--freight scheduling task as
the \emph{Selective Shared-A-Ride Pickup-and-Delivery Problem with Time Windows}
(Selective SARP-TW), a static batch-scheduling variant of the multi-compartment
vehicle routing problem with pickup, delivery and time windows in which the
planner may selectively reject orders under a limited fleet. We propose a
POMO-style multi-rollout deep reinforcement learning policy built on an
attention encoder--decoder with a shared baseline and, crucially, a
\emph{learnable reject action} that separates an explicitly priced rejection
from a silently unfulfilled order. The policy is trained under a multi-objective
reward that couples service rate, fleet size, delay and load-dependent energy,
with passengers prioritized through hard time windows. On instances of
$N\!=\!25/50/100/200$ the learned scheduler attains service rates of
$0.853/0.891/0.849/0.860$ with near-zero unfulfilled orders at small-to-medium
scale, and it exposes a controllable scale--cleanliness trade-off at
$N\!=\!200$. % TODO: refresh all numbers from the final eval2000 5-seed run.
The result is a reproducible, reject-aware DRT scheduler evaluated under a fixed
protocol together with a scale-robustness analysis.
\end{abstract}

\begin{keywords}
Demand-responsive transit; passenger--freight sharing; pickup and delivery
problem with time windows; deep reinforcement learning; attention model;
selective scheduling
\end{keywords}

% =====================================================================
\section{Introduction}\label{sec:intro}
% Funnel structure -> narrow to your gap.
% 1. Context: off-peak DRT capacity waste; urban freight pressure; SARP.
% 2. Why hard: passenger hard TW + selective acceptance + fleet minimization
%    make it a large combinatorial problem; exact/heuristic methods scale poorly.
% 3. Gap: existing DRL routing (AM, POMO) targets (C)VRP/PDP without
%    (i) passenger-priority hard TW, (ii) a *learnable* reject action, and
%    (iii) shared passenger--freight capacity in one policy.
% 4. Contributions (bullet list below).
Demand-responsive transit (DRT) has emerged as a flexible complement to fixed-route
public transport, matching vehicle supply to spatially and temporally dispersed
demand. Its efficiency, however, is highly uneven across the day. During off-peak
hours---typically the mid-day window of $10\!:\!00$ to $16\!:\!00$---passenger
demand is thin and geographically scattered, so vehicles operate with low
occupancy and a large share of capacity is wasted. In the same period, urban
logistics operators face persistent last-mile pressure and rising delivery
volumes. This asymmetry motivates \emph{passenger--freight sharing}: co-loading
parcels with passengers on the same demand-responsive fleet so that idle
off-peak capacity is monetized without procuring dedicated delivery vehicles.
The idea, often studied under the banner of the \emph{share-a-ride problem}
(SARP), promises higher asset utilization and lower operating cost, but it is
only viable if the passenger service experience is not degraded: passengers
expect to be picked up and delivered within tight, guaranteed time windows,
whereas parcels tolerate looser deadlines.

Turning this idea into a schedulable optimization problem is difficult. A
realistic off-peak co-loading operation must simultaneously respect
\emph{hard passenger time windows}, manage two physically separate
compartments (seats for passengers, cargo space for parcels), keep the
dispatched fleet small, and---because demand can exceed the profitable capacity
of a limited fleet---decide \emph{which orders to serve at all}. The last
requirement makes the problem \emph{selective}: rather than forcing every
request into the plan, the operator may reject low-value or geometrically
awkward orders to protect service quality for the rest. The resulting task is a
selective, multi-compartment pickup-and-delivery problem with time windows, a
combinatorial problem whose feasible space grows explosively with the number of
orders. Exact formulations are intractable beyond small instances, and
hand-tuned metaheuristics (simulated annealing, genetic algorithms, adaptive
large-neighborhood search) require instance-specific parameter tuning and long
run times that are ill-suited to batch replanning across many daily scenarios.

Deep reinforcement learning (DRL) offers an attractive alternative. Attention-based
neural constructive heuristics such as the Attention Model \citep{kool2018attention}
and its multi-rollout refinement POMO \citep{kwon2020pomo} learn a routing policy
once and then produce high-quality solutions for new instances in a single fast
forward pass, amortizing the optimization cost across the entire instance
distribution. These methods, however, were developed for canonical routing
problems---the (capacitated) vehicle routing problem and standard
pickup-and-delivery variants---and do not natively address three features that
are central to off-peak passenger--freight sharing: (i) \emph{passenger-priority
hard time windows} that must never be traded away for parcel throughput,
(ii) a \emph{learnable reject action} that lets the policy decline orders as a
first-class decision rather than through post-hoc filtering, and (iii) a
\emph{single shared passenger--freight capacity} model reasoned over jointly by
one policy. Closing this gap is the aim of this paper.

We formalize the off-peak shared scheduling task as \sarp{}, a static
batch-scheduling problem in which all requests are booked in advance and a
complete plan is produced before any vehicle departs, and we solve it with a
POMO-style attention policy whose decoder is augmented with a semantically
explicit reject action. The reward couples service rate, fleet size, passenger
and cargo delay, and load-dependent energy, and passengers are protected through
hard feasibility masking rather than soft penalties alone.

\paragraph{Contributions.} We make the following contributions:
\begin{itemize}
  \item We formalize \sarp{}, a static batch-scheduling variant of
        MC-VRPPDTW for off-peak shared passenger--freight DRT with hard
        passenger time windows and selective order acceptance
        (Section~\ref{sec:problem}).
  \item We design a POMO-style multi-rollout attention policy with a
        \emph{learnable reject action} (decoder logits extended by two
        semantic slots), separating an active \emph{reject} decision from
        silent \emph{unfulfilled} drops (Section~\ref{sec:method}).
  \item We define a multi-objective reward and a \emph{fixed evaluation
        protocol} (seed, greedy decoding, concurrency/dispatch limits,
        viability checks) enabling exactly reproducible benchmarking
        (Sections~\ref{sec:reward}, \ref{sec:protocol}).
  \item We report a multi-scale ($N\!=\!25/50/100/200$), multi-seed study
        and identify a controllable \emph{scale--cleanliness} trade-off,
        with comparison against SA/GA and OR baselines
        (Section~\ref{sec:exp}).
\end{itemize}

% =====================================================================
\section{Related Work}\label{sec:related}
\subsection{Shared passenger--freight transport (SARP)}
The idea of transporting people and parcels on the same vehicles was crystallized
by \citet{li2014share}, who introduced the \emph{share-a-ride problem} (SARP) as a
taxi-sharing model in which passengers and parcels are served together under time
windows, and studied a static, reduced ``freight insertion'' variant solvable by
mixed-integer programming. Because exact methods do not scale, subsequent work
turned to metaheuristics: \citet{li2016adaptive} developed an adaptive
large-neighborhood search (ALNS) for the SARP that handles substantially larger
instances, and later studies extended the model with electric and autonomous
fleets, dynamic request arrivals, and matheuristic decompositions
\citep{yu2023matheuristic}. This literature establishes co-modality as an
efficient use of transport capacity, but it predominantly targets taxi-like,
door-to-door ride-sharing and treats acceptance as exogenous. The specific
off-peak DRT setting---a small dispatched fleet, two independent
passenger/cargo compartments, and an explicit \emph{decision} of which orders to
serve---has received little attention, and almost none of it from a learning
perspective.

\subsection{Selective / prize-collecting routing with time windows}
Allowing a planner to serve only a subset of requests connects our problem to
the family of \emph{selective} or \emph{prize-collecting} routing problems. The
orienteering problem and its team and time-window extensions maximize collected
reward subject to a travel budget, and prize-collecting TSP/VRP formulations trade
served demand against routing cost \citep{vansteenwegen2011orienteering,
archetti2014routing}. These models justify selective acceptance as a principled
response to limited capacity, but they are typically single-commodity and do not
couple selection with hard passenger time windows or with a shared
passenger--freight capacity. In \sarp{} the reject decision must be made jointly
with pickup--delivery sequencing under passenger-priority windows, which is what
motivates encoding rejection directly inside the policy.

\subsection{Deep reinforcement learning for vehicle routing}
Neural constructive heuristics have become a competitive paradigm for routing.
Following pointer networks and early RL formulations, \citet{kool2018attention}
proposed the Attention Model, a Transformer-style encoder--decoder trained with
REINFORCE that solves the TSP, CVRP and several variants by autoregressive node
selection. \citet{kwon2020pomo} introduced POMO, which exploits solution symmetry
by rolling out multiple trajectories from diverse starting actions and using their
mean as a low-variance shared baseline, yielding both faster training and stronger
policies. These ideas have been extended to pickup-and-delivery and time-window
problems and, more broadly, to demand-driven mobility services
\citep{zong2021deep}. Yet existing DRL routing policies do not natively provide a
learnable order-rejection action, do not prioritize passengers through hard time
windows over co-loaded freight, and do not reason over a shared
passenger--freight capacity within a single policy. Our work fills this gap by
augmenting a POMO-style attention policy with a semantically explicit reject
action and a passenger-priority feasibility model tailored to off-peak DRT.

% =====================================================================
\section{Problem Formulation: Selective SARP-TW}\label{sec:problem}

\subsection{Setting and assumptions}\label{sec:setting}
We consider a single depot operating a homogeneous fleet of
autonomous/demand-responsive vehicles over the off-peak planning horizon
$[\underline{T},\overline{T}] = [10\!:\!00, 16\!:\!00]$. All transport
requests are booked in advance and are therefore \emph{fully known at the
start of the planning horizon}; the planner produces a complete batch
schedule before any vehicle departs. This is a \emph{static batch-scheduling}
problem: it is neither an online problem with stochastic request arrivals nor
a dynamic traffic/transit-assignment problem, and no re-optimization is
performed in response to realized events.

Each vehicle is a shared passenger--freight carrier with two independent
compartments: a passenger compartment of capacity
$Q^{\mathrm{p}}\!=\!15$ and a cargo compartment of capacity
$Q^{\mathrm{c}}\!=\!20$. Vehicles travel at constant speed
$v\!=\!25$ (distance units per hour) on a Euclidean network of side
$L\!=\!10$, incur a fixed per-node service time $s\!=\!3$~minutes, and every
trip must return to the depot within a maximum trip duration
$T^{\max}\!=\!3$~hours. We make the standard assumptions of deterministic
travel times, no intermediate depot replenishment, and instantaneous
(dis)embarking within the service time.

\subsection{Requests, network, and notation}\label{sec:notation}
The instance contains $N$ orders. Order $i$ is either a \emph{passenger}
order ($i\in\mathcal{P}$) or a \emph{parcel/cargo} order
($i\in\mathcal{C}$), with $\mathcal{P}\cup\mathcal{C}=\{1,\dots,N\}$ and a
passenger share of $60\%$. Each order $i$ is a pickup--delivery pair and is
mapped onto two graph nodes: a pickup node $i$ and a delivery node
$N\!+\!i$. Together with the depot (node $0$) this yields the node set
$\mathcal{V}=\{0,1,\dots,2N\}$, so pickups occupy indices $1,\dots,N$ and
deliveries occupy indices $N\!+\!1,\dots,2N$. Table~\ref{tab:notation}
summarizes the notation.

\begin{table}[t]
\tbl{Summary of notation.}
{\begin{tabular}{ll}\toprule
Symbol & Meaning\\ \midrule
$N$ & number of orders in the instance\\
$\mathcal{P},\mathcal{C}$ & passenger orders, cargo orders\\
$\mathcal{V}=\{0,\dots,2N\}$ & depot ($0$), pickups ($1..N$), deliveries ($N\!+\!1..2N$)\\
$d^{\mathrm{p}}_i,d^{\mathrm{c}}_i$ & passenger/cargo demand of order $i$ (signed)\\
$[e_i,l_i]$ & time window of node $i$ (early, late)\\
$Q^{\mathrm{p}},Q^{\mathrm{c}}$ & passenger ($15$) / cargo ($20$) capacity\\
$w$ & time-window width ($w\!=\!1.0$~h for both classes)\\
$T^{\max}$ & maximum trip duration ($3$~h)\\
$K_{\max}=\lceil N/6\rceil$ & fleet-size (dispatch) budget\\
$v,s$ & vehicle speed ($25$), service time ($3$~min)\\
$c^{\mathrm{v}}$ & fixed cost per dispatched vehicle ($20$)\\
$\tau^{\mathrm{ride}},\tau^{\mathrm{exc}}$ & max passenger ride ($70$~min) / excess-ride ($30$~min)\\ \bottomrule
\end{tabular}}
\label{tab:notation}
\end{table}

Passenger and cargo demands are drawn as small positive integers
(passengers mostly in $\{1,2\}$, cargo in $\{1,2,3\}$) at the pickup node
and enter with the opposite sign at the delivery node, so that a completed
order is load-neutral over its pickup--delivery pair. Time windows have a
fixed width $w\!=\!1.0$~hour; a delivery window is derived from the pickup
window plus service and travel time. Instances mix uniformly random and
clustered spatial layouts to reflect realistic off-peak demand geography.

\subsection{Decision, feasibility, and selective acceptance}\label{sec:decision}
A solution assigns each order either to a served pickup--delivery pair
embedded in some vehicle route, or to the \emph{rejected} set. A route is an
ordered sequence of node visits starting and ending at the depot; a schedule
is a set of such routes, one per dispatched vehicle. A feasible schedule must
satisfy:
\begin{enumerate}
  \item \textbf{Pairing and precedence.} For every served order, its pickup
        is visited before its delivery on the \emph{same} route.
  \item \textbf{Dual capacity.} At no point may the onboard passenger load
        exceed $Q^{\mathrm{p}}$ or the onboard cargo load exceed
        $Q^{\mathrm{c}}$; the two compartments are tracked independently.
  \item \textbf{Hard pickup time windows.} A pickup may not be served after
        its late bound $l_i$; arriving late to a pickup is infeasible for
        both passenger and cargo orders.
  \item \textbf{Passenger ride-time limits.} A passenger's in-vehicle time
        may not exceed $\tau^{\mathrm{ride}}\!=\!70$~minutes, nor exceed the
        direct-ride time by more than $\tau^{\mathrm{exc}}\!=\!30$~minutes.
  \item \textbf{Trip-duration and horizon limits.} Every trip must return to
        the depot within $T^{\max}\!=\!3$~hours, and no action may be
        scheduled to finish after the horizon end $\overline{T}\!=\!16\!:\!00$.
  \item \textbf{Fleet budget.} At most $K_{\max}=\lceil N/6\rceil$ vehicles
        (equivalently, depot dispatches) may be used.
\end{enumerate}
Delivery lateness, in contrast, is \emph{not} a hard constraint: a late
delivery is permitted but penalized (Section~\ref{sec:reward}), reflecting
the passenger-priority design in which arriving to collect a passenger on
time is mandatory while small delivery delays are tolerable.

The defining feature of the problem is \emph{selective acceptance}: because
serving every order under hard time windows and a limited fleet may be
infeasible or uneconomical, the planner is allowed to \emph{reject} a
subset of orders. Rejection is a first-class decision rather than a failure
mode, and the objective explicitly prices it (Section~\ref{sec:reward}). We
distinguish an \emph{active reject}---an order the planner deliberately
declines---from an \emph{unfulfilled} order that is left incomplete without an
explicit decision (e.g.\ because time or fleet ran out); the latter is
penalized more heavily, encouraging the policy to make its rejections
explicit and accountable.

\begin{definition}[\sarp{}]\label{def:sarp}
Given a depot, a homogeneous shared passenger--freight fleet with per-vehicle
capacities $(Q^{\mathrm{p}},Q^{\mathrm{c}})$ and dispatch budget $K_{\max}$,
and a set of $N$ advance-booked pickup--delivery orders with hard pickup time
windows, the \emph{Selective Shared-A-Ride Pickup-and-Delivery Problem with
Time Windows} asks for a batch schedule---a set of at most $K_{\max}$
depot-rooted routes together with a set of rejected orders---that satisfies
constraints (1)--(6) above and minimizes the weighted operating objective of
Section~\ref{sec:reward}, in which unserved demand and vehicle usage dominate
routing, delay, and energy costs.
\end{definition}

\sarp{} generalizes the multi-compartment vehicle routing problem with
pickup, delivery, and time windows (MC-VRPPDTW): setting
$Q^{\mathrm{c}}\!=\!0$ recovers a selective dial-a-ride problem, while
forbidding rejection recovers a standard MC-VRPPDTW. It is therefore
NP-hard, which motivates the learning-based construction heuristic developed
next.

% =====================================================================
\section{Methodology}\label{sec:method}
We solve \sarp{} by learning a construction policy that builds the schedule
one node at a time. The policy is an attention-based encoder--decoder trained
by policy gradient with a POMO-style shared baseline, and is augmented with a
learnable reject action so that selective acceptance is part of the same
sequential decision process.

\subsection{Sequential decision process}\label{sec:mdp}
We cast schedule construction as a Markov decision process solved by a single
logical vehicle that repeatedly departs from and returns to the depot; each
depot departure corresponds to one dispatched vehicle, so a full trajectory
realizes up to $K_{\max}$ routes.

\paragraph{State.} At step $t$ the state $s_t$ comprises static instance
tensors (node coordinates, node types, signed passenger/cargo demands, and
time windows) together with the dynamic quantities that evolve during
construction: the current clock $\mathrm{time}_t$, the current trip's start
time, the last visited node $a_{t-1}$, the onboard passenger and cargo loads
$(q^{\mathrm{p}}_t,q^{\mathrm{c}}_t)$, the number of vehicles already
dispatched $k_t$, a visited mask over all pickup/delivery slots, the set of
orders currently onboard (picked up but not yet delivered), per-passenger
pickup timestamps (used to enforce ride-time limits), and the rejected-order
mask with its counter. Dispatch-policy counters---the number of orders served
since the last dispatch and the number of currently open orders---are also
tracked to enforce the operational limits below.

\paragraph{Action and transition.} An action selects one graph node or the
reject slot. Selecting an order node advances the clock by the travel time
$\mathrm{dist}/v$ plus the service time $s$, updates the onboard load by the
node's signed demand, and marks the node visited. Leaving the depot fixes the
trip start time and increments the vehicle counter $k_t$; returning to the
depot resets the clock to the horizon start and clears the
served-since-dispatch counter. A reject action changes neither time,
position, nor load; it eliminates one order (Section~\ref{sec:reject}). The
episode terminates when every order is either served or rejected, the vehicle
is at the depot, and no order remains open.

\paragraph{Operational limits.} Two dispatch-shaping parameters are enforced
throughout: a minimum of $\underline{m}\!=\!4$ orders per dispatch (the
vehicle may not return to the depot before serving enough orders on the
current trip) and a cap of $\bar{o}$ concurrently open orders. These, together
with the fleet budget $K_{\max}$, keep constructed routes operationally
realistic.

\paragraph{Policy and objective.} A stochastic policy $\pi_{\bm\theta}$
parameterized by the encoder--decoder network factorizes the probability of a
complete construction trajectory
$\bm{a}=(a_1,\dots,a_T)$ over an instance $x$ autoregressively,
\begin{equation}\label{eq:policy}
\pi_{\bm\theta}(\bm{a}\mid x)
\;=\;\prod_{t=1}^{T} p_{\bm\theta}\!\left(a_t \mid s_t, x\right),
\end{equation}
where each factor $p_{\bm\theta}(a_t\mid s_t,x)$ is the masked action
distribution emitted by the decoder at step $t$ (Eq.~\ref{eq:logits}) and $T$
is the (instance-dependent) trajectory length at which the terminal condition
above is reached. Since the schedule and its rejected set are fully determined
by $\bm{a}$, the learning objective is to minimize the expected operating cost
of the induced solution,
\begin{equation}\label{eq:objective}
\min_{\bm\theta}\;\;
\mathbb{E}_{x\sim\mathcal{D}}\;
\mathbb{E}_{\bm{a}\sim\pi_{\bm\theta}(\cdot\mid x)}
\big[\,\mathrm{cost}(x,\bm{a})\,\big],
\end{equation}
with $\mathrm{cost}(\cdot)$ the weighted operating cost of
Section~\ref{sec:reward} and $\mathcal{D}$ the instance distribution of
Section~\ref{sec:data}. The gradient of Eq.~\eqref{eq:objective} is estimated
by REINFORCE with the shared baseline of Section~\ref{sec:pomo}.

\subsection{Encoder}\label{sec:encoder}
Each node is embedded from six numeric features---the two coordinates, the
passenger and cargo demands, and the early/late time-window bounds---through a
linear projection, to which a learned \emph{type embedding} (depot, passenger
pickup, cargo pickup, passenger delivery, cargo delivery) is added; the depot
uses a dedicated projection of its coordinates. The node embeddings are
processed by a graph attention encoder of $L\!=\!2$ layers, each consisting of
a multi-head self-attention sublayer ($H\!=\!8$ heads) and a feed-forward
sublayer of hidden width $512$, both wrapped in skip connections and
batch normalization, following the Transformer/Attention-Model design of
\citet{kool2018attention}.
% 待补充：核对编码器超参与代码是否一致。代码 PHASE_CONFIGS 各尺度默认
% n_encode_layers=6、hidden_dim=256、embedding_dim=256、n_heads=8;此处正文写的是
% L=2、FF=512。请确认最终上报 checkpoint 实际使用的层数/隐藏维度并统一(建议改为
% L=6、hidden=256、d=256,或注明具体尺度的取值)。 To make the encoder pickup-delivery aware, the attention
compatibility between the pickup and delivery nodes of the same order is
shifted by a learnable scalar bias, so that paired nodes attend to each other
more strongly. The encoder produces per-node embeddings and a graph embedding
(their mean) and is run \emph{once} per instance.

Concretely, each node $i$ is first mapped to an initial embedding
$\bm{h}_i^{(0)}\in\mathbb{R}^{d}$ ($d$ the embedding dimension). For an order
node this is the sum of a linear projection of its five numeric features
$\bm{f}_i=(x_i,y_i,d^{\mathrm{p}}_i,d^{\mathrm{c}}_i,e_i,l_i)$ and a learned
type embedding $\bm{g}_{\tau(i)}$ selected from the five node types
$\tau\in\{\text{depot},\text{pass.\ pickup},\text{cargo pickup},\text{pass.\
delivery},\text{cargo delivery}\}$; the depot uses a dedicated coordinate
projection plus its type embedding:
\begin{equation}\label{eq:init-embed}
\bm{h}_i^{(0)} = \bm{W}^{\mathrm{num}}\bm{f}_i + \bm{g}_{\tau(i)},
\qquad
\bm{h}_0^{(0)} = \bm{W}^{\mathrm{dep}}(x_0,y_0)^{\!\top} + \bm{g}_{\text{depot}}.
\end{equation}
The embeddings are refined by $L$ Transformer layers. Layer $\ell$ applies a
multi-head self-attention (MHA) sublayer and a node-wise feed-forward (FF)
sublayer, each wrapped in a skip connection and batch normalization (BN):
\begin{align}\label{eq:enc-layer}
\hat{\bm{h}}^{(\ell)} &= \mathrm{BN}\!\big(\bm{h}^{(\ell-1)} + \mathrm{MHA}^{(\ell)}(\bm{h}^{(\ell-1)})\big),\\
\bm{h}^{(\ell)} &= \mathrm{BN}\!\big(\hat{\bm{h}}^{(\ell)} + \mathrm{FF}^{(\ell)}(\hat{\bm{h}}^{(\ell)})\big).
\end{align}
Within each attention head $m$, queries, keys, and values are linear
projections of the layer input, and the pickup--delivery bias $\beta$ is added
to the compatibility of nodes belonging to the same order before the softmax,
\begin{equation}\label{eq:enc-attn}
u^{m}_{ij} = \frac{(\bm{W}^{Q}_m\bm{h}_i)^{\!\top}(\bm{W}^{K}_m\bm{h}_j)}{\sqrt{d_k}}
+ \beta\,\mathbb{1}[\,i,j\ \text{are a pickup--delivery pair}\,],
\qquad
\bm{a}^{m}_{i} = \mathrm{softmax}_j\big(u^{m}_{ij}\big),
\end{equation}
with $d_k=d/H$ the per-head key dimension and $H$ the number of heads. The
final node embeddings $\bm{h}_i=\bm{h}_i^{(L)}$ and the graph embedding
$\bar{\bm{h}}=\tfrac{1}{|\mathcal{V}|}\sum_i \bm{h}_i$ are passed to the
decoder.

\subsection{Decoder and context}\label{sec:decoder}
The decoder builds the schedule autoregressively. At each step it forms a
context vector by concatenating the embedding of the current node with four
dynamic scalars: the remaining passenger capacity $Q^{\mathrm{p}}\!-\!q^{\mathrm{p}}_t$,
the remaining cargo capacity $Q^{\mathrm{c}}\!-\!q^{\mathrm{c}}_t$, the
normalized clock, and the remaining vehicle budget
$(K_{\max}\!-\!k_t)/K_{\max}$. A glimpse-style multi-head attention over the
node embeddings yields the node compatibility logits; these are $\tanh$-clipped
with clipping constant $C\!=\!10$ and combined with the feasibility mask of
Section~\ref{sec:masking}. The next node is chosen greedily (arg-max) at
evaluation or by sampling during training.

Formally, the step context $\bm{c}_t$ concatenates the embedding of the
current node $\bm{h}_{a_{t-1}}$ with the four dynamic scalars, and the decoder
query is
\begin{equation}\label{eq:dec-query}
\bm{c}_t = \big[\,\bm{h}_{a_{t-1}} \,\big\Vert\,
Q^{\mathrm{p}}\!-\!q^{\mathrm{p}}_t,\;
Q^{\mathrm{c}}\!-\!q^{\mathrm{c}}_t,\;
\widetilde{\mathrm{time}}_t,\;
(K_{\max}\!-\!k_t)/K_{\max}\,\big],
\qquad
\bm{q}_t = \bm{W}^{C}\bar{\bm{h}} + \bm{W}^{S}\bm{c}_t,
\end{equation}
where $\widetilde{\mathrm{time}}_t\!\in\![0,1]$ is the clock normalized over
the operating horizon and $\bm{W}^{C},\bm{W}^{S}$ are learned projections
(the fixed-context term $\bm{W}^{C}\bar{\bm{h}}$ is precomputed once per
instance). A multi-head glimpse then attends over the node embeddings to
refine the query into $\tilde{\bm{q}}_t$, and single-head compatibilities with
the node keys $\bm{k}_j=\bm{W}^{K}\bm{h}_j$ give the per-node logits
\begin{equation}\label{eq:dec-logits}
z^{\mathrm{node}}_{t,j}
= \frac{\tilde{\bm{q}}_t^{\!\top}\bm{k}_j}{\sqrt{d}}\,,
\qquad j\in\{0,1,\dots,2N\}.
\end{equation}
These node logits are concatenated with the reject logit of
Section~\ref{sec:reject}, clipped, masked (Eq.~\ref{eq:logits}), and passed
through a softmax to give the action distribution
\begin{equation}\label{eq:dec-softmax}
p_{\bm\theta}(a_t = j \mid s_t, x)
= \frac{\exp(z_{t,j})}{\sum_{j'} \exp(z_{t,j'})}\,,
\end{equation}
which is the factor appearing in Eq.~\eqref{eq:policy}.

\subsection{Learnable reject action}\label{sec:reject}
Selective acceptance is realized by appending a single extra logit to the
$2N\!+\!1$ node logits (depot plus $N$ pickups and $N$ deliveries), yielding a
$(2N\!+\!2)$-way output. The extra logit is produced by a dedicated linear
head applied to the decoder context, whose bias is initialized negatively
($-2.5$) so that rejection is rare early in training and is learned only when
it is beneficial. Formally,
\begin{equation}\label{eq:logits}
\bm{z} \;=\; \big[\,\bm{z}^{\mathrm{node}}\,\big\Vert\, \phi_{\mathrm{rej}}(\bm{c}_t)\,\big],
\qquad
\bm{z} \leftarrow C\cdot\tanh(\bm{z}),
\qquad
z_j \leftarrow -\infty \;\; \forall j\in\mathcal{M}_t,
\end{equation}
where $\bm{c}_t$ is the decoder context, $\phi_{\mathrm{rej}}$ the reject head,
and $\mathcal{M}_t$ the set of masked (infeasible) actions.

When the reject slot is selected, the \emph{target} order is chosen
deterministically as the untouched, not-yet-rejected, not-onboard pickup with
the earliest late time-window bound---i.e.\ the most time-critical
outstanding order---so the policy learns \emph{when} to reject while the
resolution of \emph{which} order is fixed and interpretable. Denoting by
$\mathcal{R}_t$ the set of orders whose pickup is unvisited, not onboard, and
not already rejected at step $t$, the rejected order is
\begin{equation}\label{eq:reject-target}
i^{\star}_t \;=\; \arg\min_{i\,\in\,\mathcal{R}_t}\; l_i,
\end{equation}
with $l_i$ the late bound of pickup $i$; if $\mathcal{R}_t=\varnothing$ the
reject slot is masked out. Rejection is
offered by the mask only when a reject candidate exists and the vehicle is
either pre-departure or in a dead-end, preventing rejections in the middle of
a committed trip.

This design cleanly separates two failure semantics that are conflated in
naive formulations. An \emph{active reject} is an order eliminated through the
reject action and priced at $\alpha_{R}$; an \emph{unfulfilled} order is one
left incomplete without an explicit reject and priced at the strictly larger
$\alpha_{U}\!>\!\alpha_{R}$ (Section~\ref{sec:reward}). Because silent failure
costs more than an explicit decline, the policy is driven to convert
would-be unfulfilled demand into accountable rejections.

\subsection{Feasibility masking}\label{sec:masking}
All hard constraints of Section~\ref{sec:decision} are enforced exactly at
decoding time through the mask $\mathcal{M}_t$, so the policy can only ever
emit feasible actions. The mask forbids: already-visited nodes; deliveries
whose pickup is not yet done and pickups of rejected orders (precedence);
pickups whose demand exceeds the remaining passenger or cargo capacity;
pickups reached after their late bound (hard time window); actions that would
violate the passenger ride-time or excess-ride-time limits; actions whose
predicted trip duration exceeds $T^{\max}$ or whose predicted finish exceeds
the horizon end; depot returns while the vehicle is carrying orders or before
the minimum-orders-per-dispatch threshold is met; and all order nodes once the
fleet budget $K_{\max}$ is exhausted. A bitmask feasibility-completion check
additionally masks pickups whose committed delivery could not be completed
within the open-order and time limits. Delivery lateness is deliberately left
unmasked and handled as a soft penalty.

\subsection{POMO-style training with a shared baseline}\label{sec:pomo}
We train with REINFORCE using a POMO-style shared baseline. The encoder is run
once per instance and its embeddings are replicated into multiple parallel
decoding rollouts (multi-start decoding rather than input augmentation). The
baseline for each instance is the mean return over its rollouts, so the
advantage of a rollout is its return minus this shared mean; this removes the
need for a separate critic network and reduces gradient variance. The policy
parameters are updated by the standard policy-gradient estimator over the
batch. Algorithm~\ref{alg:train} summarizes one training epoch.

Concretely, for instance $x$ with $R$ rollouts $\bm{a}^{(1)},\dots,\bm{a}^{(R)}$
sampled from $\pi_{\bm\theta}$, let $\mathrm{cost}_r=\mathrm{cost}(x,\bm{a}^{(r)})$.
The shared baseline is the mean cost over the rollouts and the gradient
estimator is
\begin{equation}\label{eq:pg}
b(x)=\frac{1}{R}\sum_{r=1}^{R}\mathrm{cost}_r,
\qquad
\nabla_{\bm\theta}\mathcal{L}
=\mathbb{E}_{x}\!\left[\frac{1}{R}\sum_{r=1}^{R}
\big(\mathrm{cost}_r - b(x)\big)\,
\nabla_{\bm\theta}\log\pi_{\bm\theta}(\bm{a}^{(r)}\mid x)\right],
\end{equation}
where $\log\pi_{\bm\theta}(\bm{a}^{(r)}\mid x)=\sum_{t}\log p_{\bm\theta}(a^{(r)}_t\mid s_t,x)$
is the trajectory log-likelihood (Eq.~\ref{eq:policy}). When only a single
rollout per instance is used ($R\!=\!1$), the baseline reduces to the
batch-mean cost, i.e.\ $b=\tfrac{1}{|\mathcal{B}|}\sum_{x\in\mathcal{B}}\mathrm{cost}(x)$,
which is the configuration used for the runs reported in this paper.
% 待补充：确认最终上报结果所用的 pomo_size / baseline_mode（代码默认 pomo_size=1、
% baseline_mode=auto → batch-mean 基线）。若正文其他处已声明使用 per-instance
% multi-rollout 共享基线，请与此处统一：目前代码 PHASE_CONFIGS 各尺度均为 pomo_size=1。

\begin{algorithm}[t]
\caption{Reject-aware POMO-style training (one epoch)}\label{alg:train}
\begin{algorithmic}[1]
\State \textbf{Input:} instance batch $\mathcal{B}$, rollouts per instance $R$, policy $\pi_\theta$
\For{each instance $x\in\mathcal{B}$}
  \State encode $x$ once $\rightarrow$ node/graph embeddings
  \State replicate embeddings into $R$ parallel rollouts
  \For{$r=1,\dots,R$}
    \State construct schedule by autoregressive decoding with mask $\mathcal{M}_t$ (Eq.~\ref{eq:logits})
    \State record return $G_r = -\,\mathrm{cost}(x,\text{schedule}_r)$ (Eq.~\ref{eq:reward})
  \EndFor
  \State shared baseline $b(x) \gets \tfrac{1}{R}\sum_{r} G_r$
  \State accumulate $\nabla_\theta \gets \sum_{r}(G_r - b(x))\,\nabla_\theta \log \pi_\theta(\text{schedule}_r\mid x)$
\EndFor
\State update $\theta$ by gradient ascent on the accumulated estimator
\end{algorithmic}
\end{algorithm}

\begin{table}[t]
\tbl{Network and optimization hyperparameters. Values are read from the
per-scale training configuration; unless overridden on the command line, all
scales share the architecture and differ only in the epoch/batch schedule.}
{\begin{tabular}{lc}\toprule
Hyperparameter & Value \\ \midrule
Embedding dimension $d$        & $256$ \\
Attention heads $H$            & $8$ \\
Encoder layers $L$             & $6$ \\
Feed-forward hidden width      & $256$ \\
Normalization                  & batch \\
$\tanh$-clipping $C$           & $10$ \\
Reject-head initial bias       & $-2.5$ \\
Rollouts per instance $R$ (\texttt{pomo\_size}) & $1$ \\
Baseline mode                  & batch-mean (\texttt{auto}) \\
Optimizer                      & AdamW \\
Learning-rate schedule         & cosine annealing \\
Initial learning rate          & $5\times10^{-5}$ \\ \bottomrule
\end{tabular}}
\label{tab:hparams}
\end{table}
% 待补充：核对表中 hidden_dim / n_encode_layers / lr 是否随尺度变化。代码
% PHASE_CONFIGS[25] 给出 lr=5e-5、hidden_dim=256、n_encode_layers=6、batch_size=8、
% n_epochs=50、epoch_size=8000;优化器为 AdamW + CosineAnnealingLR。若 N=50/100/200 的
% lr、batch_size、n_epochs 与之不同,请补出各尺度取值或在表注中说明。

% =====================================================================
\section{Reward Design}\label{sec:reward}
The episode return is the negative of a weighted operating cost that
aggregates six terms: energy, delay, vehicle usage, active rejection,
unfulfilled demand, and trip overtime. Writing $C_E$ for energy,
$C_D$ for delay, $K$ for the number of dispatched vehicles,
$n_{\mathrm{rej}}$ and $n_{\mathrm{unf}}$ for the active-reject and
unfulfilled counts, and $C_T$ for aggregate trip overtime, the cost is
\begin{equation}\label{eq:reward}
\mathrm{cost} \;=\;
  \alpha_{E}\,\tilde{C}_{E}
+ \alpha_{D}\,\big(\tilde{C}_{D}^{\mathrm{p}}+\tilde{C}_{D}^{\mathrm{c}}\big)
+ \alpha_{V}\,\tilde{K}
+ \alpha_{R}\,n_{\mathrm{rej}}
+ \alpha_{U}\,n_{\mathrm{unf}}
+ \alpha_{T}\,C_{T},
\qquad R=-\,\mathrm{cost},
\end{equation}
where $\tilde{(\cdot)}$ denotes size-dependent normalization of the routing,
delay, and vehicle terms (Section~\ref{sec:normalization}) and the count-based
penalties are applied at their raw scale. Table~\ref{tab:reward} lists the
coefficients and the physical configuration; all experiments in this paper use
these values.

\begin{table}[t]
\tbl{Reward coefficients and physical configuration (identical across all runs).}
{\begin{tabular}{lcccccccc}\toprule
$\alpha_E$ & $\alpha_D$ & $\alpha_V$ & $\alpha_R$ & $\alpha_U$ & $\alpha_T$ & $Q^{\mathrm{p}}$ & $Q^{\mathrm{c}}$ & $c^{\mathrm{v}}$\\ \midrule
$1.0$ & $2.5$ & $3.0$ & $575$ & $750$ & $200$ & $15$ & $20$ & $20$\\ \bottomrule
\end{tabular}}
\label{tab:reward}
\end{table}

\subsection{Term definitions}\label{sec:terms}
\paragraph{Energy $C_E$.} Energy is load-dependent. Along each arc the onboard
weight combines passengers (at $65$~kg each) and parcels (at unit weight),
and the per-distance energy rate scales with this weight,
$\rho = \rho_0\,(1 + m_{\mathrm{kg}}/m_{\mathrm{ref}})$ with base rate
$\rho_0\!=\!0.18$ and reference mass $m_{\mathrm{ref}}\!=\!10^4$. The raw
energy cost is $\sum_{\text{arcs}} \rho\cdot \mathrm{dist}$ times the unit
electricity price; cumulative load resets to zero at each depot visit.

\paragraph{Delay $C_D$.} Only \emph{deliveries} accrue delay, computed as the
positive part of arrival time minus the delivery's late bound. Passenger
delivery delay is priced linearly at $0.6$ per minute. Cargo delivery delay is
priced by a three-tier convex schedule---$0.06$ per minute up to $30$~min,
$0.18$ per minute from $30$ to $60$~min, and $0.36$ per minute beyond
$60$~min---so that large parcel delays are penalized progressively while small
ones remain cheap. Pickup lateness does not appear here: it is a hard
constraint enforced by masking (Section~\ref{sec:masking}).

\paragraph{Vehicle usage $K$.} Each pickup that is immediately preceded by a
depot visit counts as a new dispatch; the raw vehicle cost is
$K\cdot c^{\mathrm{v}}$ with $c^{\mathrm{v}}\!=\!20$. Because $\alpha_V$
multiplies the normalized count, the objective actively pressures the policy
to consolidate demand into fewer trips.

\paragraph{Trip overtime $C_T$.} At every depot return and at the final step,
the positive part of the trip duration minus $T^{\max}\!=\!3$~h is accumulated;
its weight $\alpha_T\!=\!200$ is large, effectively discouraging any overtime
while keeping the constraint differentiable through the reward rather than
purely through masking.

\paragraph{Rejection and unfulfilled demand.} Let $u$ be the number of
unserved orders. The active-reject count is
$n_{\mathrm{rej}}=\min(\text{reject\_count},\,u)$ and the unfulfilled count is
$n_{\mathrm{unf}}=\max(u-n_{\mathrm{rej}},\,0)$. These are priced at
$\alpha_R\!=\!575$ and $\alpha_U\!=\!750$ respectively. The strict ordering
$\alpha_U>\alpha_R\gg$ (routing/delay costs) encodes the design priority:
serving demand dominates operational efficiency, and among unserved orders an
explicit, accountable rejection is preferred to a silent unfulfilled drop.

\subsection{Normalization}\label{sec:normalization}
To make returns comparable across instance sizes during training, the routing,
delay, and vehicle terms are divided by size-dependent normalization profiles
keyed on the graph size ($N\!\in\!\{25,50,100\}$), while the count-based
reject/unfulfilled/overtime penalties are kept at their raw scale so that the
relative priority in Eq.~\eqref{eq:reward} is preserved regardless of $N$. We
report both the normalized training objective and the raw (monetary) cost
decomposition in the experiments.

% =====================================================================
\section{Experimental Setup}\label{sec:setup}

\subsection{Instances and data}\label{sec:data}
We evaluate on synthetic instances at four scales, $N\in\{25,50,100,200\}$.
Each instance places a depot and $N$ pickup--delivery orders on a
$10\times10$ Euclidean area, with $60\%$ passenger orders and $40\%$ cargo
orders. Spatial layouts mix uniformly random and clustered demand (three
clusters) in equal proportion to reflect off-peak geography. Passenger demands
are small integers (mostly $\{1,2\}$), cargo demands lie in $\{1,2,3\}$, and
each order carries a hard pickup time window of width $w\!=\!1.0$~h sampled
under a period-mixed scheme over the $10\!:\!00$--$16\!:\!00$ horizon;
delivery windows are derived from the pickup window plus service and travel
time. For external comparability we additionally map standard Li~\&~Lim and
Cordeau pickup-delivery benchmarks onto our setting, and use a cleaned
real-world demand trace only as an illustrative case study rather than as the
primary benchmark, so that reported performance is not tied to a single
proprietary dataset.

\subsection{Fixed evaluation protocol}\label{sec:protocol}
To make results exactly reproducible and cross-comparable, every evaluation in
this paper uses one frozen protocol: fixed evaluation seed $99999$; greedy
(arg-max) decoding; at most $\bar{o}\!=\!6$ concurrently open orders; a
minimum of $\underline{m}\!=\!4$ orders per dispatch; delivery-viability
checking enabled; and the viability fallback disabled. We report on two
evaluation-set sizes, a fine-grained set of $500$ instances and a primary set
of $2000$ instances, and average over five seeds for stability where noted.

A schedule is deemed \emph{clean} only if it simultaneously passes four gates:
zero unfulfilled orders, zero pickup-only (picked-up-but-never-delivered)
orders, zero started-but-not-completed orders, and zero
untouched-yet-unrejected orders. This all-or-nothing cleanliness criterion is
stricter than service rate alone and is used to diagnose the semantic quality
of a schedule, not merely how many orders it happened to serve.

\subsection{Baselines and ablations}\label{sec:baselines}
We compare against two metaheuristics---simulated annealing (SA) and a genetic
algorithm (GA)---configured for the same objective and constraints, and, where
tractable, an exact/OR reference on small instances. To isolate the
contribution of each design choice we run three ablations: (i) removing the
learnable reject action; (ii) removing the shared baseline (single-rollout
REINFORCE); and (iii) varying the number of rollouts per instance. To
pre-empt any concern about checkpoint cherry-picking, we additionally report,
under the fixed protocol, the service rate of the service-selected, final,
objective-selected, and best-loss checkpoints, showing that the reported
conclusions are not an artifact of checkpoint selection.

\subsection{Metrics}\label{sec:metrics}
We report: \emph{service rate} (fraction of orders served); the number of
dispatched vehicles $K$; mean passenger and cargo delivery delay; the
unfulfilled-order ratio; and the clean-gate pass rate defined above. Together
these capture the three axes the objective trades off---demand coverage, fleet
economy, and service quality---and expose the scale--cleanliness tension
analyzed in Section~\ref{sec:exp}.

% =====================================================================
\section{Results and Discussion}\label{sec:exp}
\subsection{Main results across scales}
% eval2000, 5 seeds per scale (99999,10007,10037,10067,10099);
% replay audit partition-consistent 2000/2000 on every seed.
\begin{table}[t]
\tbl{Main results of the proposed method across instance scales, under the
fixed evaluation protocol of Sec.~\ref{sec:protocol}. Each row is the mean over
five seeds ($\{99999,10007,10037,10067,10099\}$), each evaluated on $2000$
greedy-decoded instances; the semantic replay audit is partition-consistent on
$2000/2000$ instances for every seed. ``Clean'' denotes that all four
business-cleanliness gates pass.}
{\begin{tabular}{lccccc}\toprule
Scale $N$ & Service rate & Reject rate & Unfulfilled & \#Vehicles & Clean \\ \midrule
25  & $0.853\pm0.001$ & $0.147$ & $0.0000$ & $4.19$  & yes         \\
50  & $0.889\pm0.001$ & $0.112$ & $0.0000$ & $7.96$  & yes         \\
100 & $0.844\pm0.003$ & $0.156$ & $0.0004$ & $14.84$ & borderline  \\
200 & $0.859\pm0.004$ & $0.114$ & $0.0272$ & $27.55$ & no (mild)   \\ \bottomrule
\end{tabular}}
\label{tab:main}
\end{table}

\subsection{Scale--cleanliness trade-off}
% N25/N50 clean & high-service -> N100 borderline -> N200 mild controllable
% degradation. Note: at N200 pickup_only=0 and started_not_completed=0
% (no severe semantic breakdown); violations confined to
% unfulfilled/untouched_unrejected (~5.5/200 orders, same batch);
% passenger hard-TW & ride-time violations = 0. Frame as a *finding*.

\subsection{Comparison with baselines}
% DRL (5-seed eval2000 mean) vs SA/GA under the SAME reject-enabled protocol.
% SA/GA: 500 aligned instances, reject action enabled (comparable to DRL).
% N200 SA/GA aligned-500 runs pending -> placeholders below.
We compare the proposed method against two metaheuristic baselines, simulated
annealing (SA) and a genetic algorithm (GA), under an identical
reject-enabled protocol so that all three methods may decline orders rather
than being forced to serve every request. For the learned policy we report the
five-seed \texttt{eval2000} mean; the metaheuristics are evaluated on $500$
aligned instances. Table~\ref{tab:baselines} summarizes the comparison.

\begin{table}[t]
\tbl{Comparison of the proposed method (Ours) with simulated annealing (SA)
and a genetic algorithm (GA) under the shared reject-enabled protocol. Lower
raw cost and fewer vehicles are better; higher service rate is better.
``Unf.'' is the residual unfulfilled rate. Ours: five-seed \texttt{eval2000}
mean. SA/GA: $500$ aligned instances. Results for $N{=}200$ SA/GA are pending
and marked ``--''.}
{\begin{tabular}{llccccc}\toprule
Scale $N$ & Method & Service & Reject & Unf. & \#Veh. & Raw cost \\ \midrule
\multirow{3}{*}{25}
 & Ours & $\mathbf{0.853}$ & $0.147$ & $0.0000$ & $\mathbf{4.19}$  & $2256.3$ \\
 & SA   & $0.895$          & $0.105$ & $0.0000$ & $4.69$           & $\mathbf{1645.1}$ \\
 & GA   & $0.719$          & $0.000$ & $0.2806$ & $5.00$           & $5388.9$ \\ \midrule
\multirow{3}{*}{50}
 & Ours & $\mathbf{0.889}$ & $0.112$ & $0.0000$ & $\mathbf{7.96}$  & $\mathbf{3463.1}$ \\
 & SA   & $0.868$          & $0.131$ & $0.0011$ & $8.45$           & $4037.9$ \\
 & GA   & $0.727$          & $0.000$ & $0.2729$ & $9.00$           & $10469.9$ \\ \midrule
\multirow{3}{*}{100}
 & Ours & $\mathbf{0.844}$ & $0.156$ & $0.0004$ & $\mathbf{14.84}$ & $\mathbf{9734.1}$ \\
 & SA   & $0.719$          & $0.281$ & $0.0000$ & $16.92$          & $16626.8$ \\
 & GA   & $0.719$          & $0.000$ & $0.2811$ & $17.00$          & $21536.6$ \\ \midrule
\multirow{3}{*}{200}
 & Ours & $0.859$ & $0.114$ & $0.0272$ & $27.55$ & $18339.8$ \\
 & SA   & --      & --      & --       & --      & --        \\
 & GA   & --      & --      & --       & --      & --        \\ \bottomrule
\end{tabular}}
\label{tab:baselines}
\end{table}

Two consistent patterns emerge. First, the genetic algorithm fails to produce a
serviceable schedule at any scale: its service rate stays near $0.72$ and, unlike
the learned policy, the unserved mass leaks into residual \emph{unfulfilled}
orders ($\approx 0.28$) rather than explicit pre-service rejections, while its
raw cost is roughly $2$--$2.4\times$ that of the proposed method. Second,
simulated annealing is a genuine competitor only at the smallest scale: at
$N{=}25$ it attains a slightly higher service rate ($0.895$ vs.\ $0.853$) at the
expense of more vehicles ($4.69$ vs.\ $4.19$), but from $N{=}50$ onward the
learned policy overtakes it on service rate, vehicle usage, and cost
simultaneously, and the gap widens sharply at $N{=}100$, where SA collapses to
$0.719$ service while using more vehicles and about $1.7\times$ the cost. This
supports our central claim that the learned reject-aware policy scales more
gracefully than metaheuristic search on the selective passenger--freight
problem. % TODO: extend narrative once N200 SA/GA numbers are available.

\subsection{Ablation studies}
% learnable reject; shared baseline; #rollouts.  [TODO]

\subsection{Checkpoint-selection robustness}
% Report model_best_service vs model_best vs model_final vs
% model_best_objective at eval500 to pre-empt "cherry-picked checkpoint".

% =====================================================================
\section{Conclusion}\label{sec:conclusion}
% Recap contributions; state the scale-cleanliness limitation honestly;
% future work: dynamic/online arrivals, larger fleets, real-network case.
\emph{[TODO.]}

% ---- unnumbered back-matter sections (de-anonymized version only) ----
\section*{Acknowledgements}
% [TODO]

\section*{Disclosure statement}
No potential conflict of interest was reported by the author(s).

\section*{Funding}
% [TODO if applicable]

% =====================================================================
\bibliographystyle{tfnlm}
\bibliography{refs}   % your .bib file; start from interactnlmsample.bib

\end{document}